\documentclass[11pt]{article}
\usepackage[margin=1.1in]{geometry}
\usepackage{microtype}
\usepackage{hyperref}
\usepackage{parskip}
\usepackage{amsmath}

\title{}
\author{}
\date{}

\begin{document}

\noindent\hfill May 2026

\bigskip
\noindent
Seth,

\medskip

I've been going back through old files this spring and one of the
things I picked up again is the trademark-novelty index we sketched
out years ago and never wrote up --- the
\texttt{ericsilver/novelty} repo on GitHub, which until a few weeks
ago contained nothing but a 22~kB README that read like an unfunded
research proposal. I've now actually built the thing, and I think
the surprising part is not the part I expected to be surprising. I
wanted to write to you about it before deciding whether to push
toward a paper.

\subsubsection*{What I went looking for}

The original idea was the one we talked about: \emph{operationalize
firm innovation by reading the text trademark applicants are
forced to write about themselves}. Trademarks are unusually
well-incentivized text --- the registrant wants the goods/services
description to be broad enough to defend against future use but
narrow enough to clear the examiner --- so the corpus is a clean,
uniform, public-domain signal of what each firm claims to be at
filing time. The hypothesis I started with was ``firms that file
unusually-novel goods/services descriptions are firms that go on
to be successful.'' I was expecting either a positive linear
relationship between novelty and various firm outcomes, or, if
the pioneer-vs-fast-follower folk wisdom held, an \emph{inverse}
linear relationship: pioneers pay the cost, fast followers reap the
gain.

I set up the pipeline: USPTO Open Data Portal (16.5~M trademark
filings, 5.7~M unique registrants, all 45 NICE international
classes, 1985--2026), a per-class unigram + bigram dictionary, and
a Kullback--Leibler surprise score for every filing computed
against a 5-year look-behind reference distribution
(\emph{prospective} surprise: how unusual was the wording at the
time?) and a 5-year look-ahead reference (\emph{retrospective}
surprise: how unusual is the wording in hindsight, after the
industry has had five years to catch up?). The composite score
$\Delta KL = KL_{\text{pros}} - KL_{\text{retr}}$ is positive when
the future starts to look like the filing --- when a novel filing
gets imitated --- and negative when the filing was conservative at
the time and stayed conservative. I then linked the matched
public-firm subset (about 14{,}000 USPTO owners, 1{,}600 with
clean SEC reporting) to EDGAR financials and to monthly stock
prices.

\subsubsection*{What I found that I didn't expect}

The linear story is wrong. The relationship between $\Delta KL$
and survival outcomes is not monotonic; it is a clean
\textbf{U-shape}, present in 38 of the 44 industries with enough
data to estimate it. Both ends of the dial outperform the middle.
A trademark filing in the bottom decile of $\Delta KL$
(maximally-conservative; \texttt{Beer} on a beer; \texttt{still
water; sparkling water} on a water) clears registration at
$48.4\%$ and is renewed at ten years $35.7\%$ of the time. A
filing in the top decile (maximally-imitated-novel) clears at
$55.5\%$ and is renewed at $41.1\%$. Filings in deciles four
through six --- neither novel enough to claim a category nor
templated enough to coast on legal-clarity grounds --- do worst,
$39\%$ for both outcomes. Pooled across 8.6~M filings, the
quadratic coefficient on $z(\Delta KL)^2$ is $+0.009$ (SE
$0.0002$); per-industry, $86\%$ of industries replicate the sign,
$66\%$ at $p<0.05$. A two-sided exponential fit
$\log(\text{rate}) = a + b\,|\Delta KL|$ on the bin-aggregated data
gives positive slopes on both halves, $b_{\text{neg}}=+0.04$ and
$b_{\text{pos}}=+0.05$, both $R^2>0.7$.

The financial pattern follows the same shape but it is dominated
by the upper arm. On the matched panel of $1{,}597$ public firms,
a one-standard-deviation increase in three-year-trailing $\Delta
KL$ is associated with $+3.2$ percentage points of gross margin
(SE $0.6$); on the firm-debut variant ---restricting to each
firm's first-ever filing so that current novelty cannot have been
funded by the proceeds of past novelty --- the effect attenuates
to $+1.7$ pp but stays significant. \emph{Stock returns are where
it gets interesting.} Per $\sigma$ of debut $\Delta KL$, excess
return over the S\&P~500 grows from $+3.6$ pp at one year to
$+19$ pp at three to $+28$ pp at four, then mean-reverts to
$+15$ pp at five. The peak is several years out, not immediate;
the $R^2$ contribution is small but real ($+0.3$ pp); a separate
within-firm variance regression confirms the symmetric prediction
that high-$\Delta KL$ firms have wider gross-margin spread, i.e.\
$\Delta KL$ measures something that raises both the mean and the
variance of outcomes.

\subsubsection*{What the construct actually is}

I started this expecting to validate ``trademark novelty as a
proxy for innovation.'' I think the construct is something
slightly different and more interesting.

It is essentially uncorrelated with the dominant existing
innovation proxy --- patent counts ($r=+0.01$ on
$174{,}569$ matched firm-years from PatentsView) --- and it is
near-zero on Crunchbase total funding, weakly positive on
R\&D-to-revenue ($r=+0.13$), and modestly elevated on the MIT
Technology Review 50 Smartest Companies list ($+0.07$ on-list vs.\
off-list). Yet on stock returns, an $R^2$ decomposition shows
patents and $\Delta KL$ are essentially additive --- their joint
contribution to $R^2$ is the sum of the two singletons --- and
prospective surprise alone (mere novelty without the
later-imitation marker) adds nothing whatsoever. The forward-looking
``the-future-came-to-resemble-this'' component is what makes
$\Delta KL$ predictive.

Read together with the U-shape, the cleanest summary I have is
that $\Delta KL$ measures \textbf{the smartness, and the
riskiness, of the bets a firm makes when it names its products}.
Average bets earn average outcomes. Bets that depart from the
industry-mean vocabulary in either direction earn larger or
smaller outcomes than average; the direction depends on whether
the rest of the industry comes to imitate the bet, but the
magnitude depends only on how far the firm departed from the
centre. The five-year half-life on returns suggests the rewards
take time to ferment because imitators have to recognise the new
category and copy it before the originator's claim becomes
financially valuable.

This is, I think, an evidentiary claim that ideas matter,
literally. A firm whose trademark filings sit at the centre of
the industry vocabulary distribution earns market-average margins
and market-average returns. A firm whose filings depart from the
centre, in either direction, has measurably different financial
dynamics. That this shows up across 45 NICE classes and 1{,}500
matched public firms suggests it is a population-level pattern,
not a quirk of software-and-pharma.

\subsubsection*{Should this be a paper?}

This is the question I want to ask you.

The honest framing of the contribution, as I see it now, is:

\begin{enumerate}
  \item A new \textbf{publicly-replicable, complete-coverage,
        real-time} firm-level innovation measure derived from
        cheap textual data with strong economic incentives behind
        it. The pipeline is in a public GitHub repo with a
        Makefile that reproduces every figure and table from the
        raw USPTO and SEC sources.
  \item Empirical demonstration that the relationship between
        novelty and firm outcomes is \textbf{not} linear, nor a
        clean ``pioneers vs.\ fast followers'' tradeoff, but a
        U-shape: \emph{the middle of the distribution is the
        worst place to be}. To my knowledge no prior work in the
        innovation-measurement literature has documented this.
  \item Demonstration that the new measure is \textbf{additive}
        with patent counts in predicting stock returns; the two
        capture different innovation channels.
  \item Demonstration that the predictive horizon is several
        years, peaking at four, with mean reversion at five;
        consistent with a real category-creation/imitation cycle
        rather than a contemporaneous correlation artifact.
  \item A reusable framework: the same Kullback--Leibler
        prospective/retrospective decomposition can be applied to
        any other dated text corpus in which the firm or
        individual is incentivised to be precise --- 10-K business
        descriptions, app-store descriptions, NIH grant abstracts,
        Kickstarter project descriptions.
\end{enumerate}

The honest concerns:

\begin{enumerate}
  \item The construct is a \emph{language} novelty measure, not a
        product or technology novelty measure. A firm that solves
        an old problem with an entirely new device but describes
        it in old vocabulary is invisible. Liquid Death's canned
        water is a clean example --- the goods/services text is
        ``still water; sparkling water'' and $\Delta KL$ is
        almost zero, even though the brand was a category
        creator.
  \item The matched-firm panel ($1{,}600$ firms) is small relative
        to the SEC universe ($\sim 13{,}000$ firms), and the
        match itself is fully automated --- normalised exact
        match plus $\texttt{rapidfuzz}$ token-sort with cutoff
        92, no manual review. There are inevitable false
        positives and false negatives.
  \item The $R^2$ contributions to stock returns are real but
        small (under $0.5$ pp incremental at any horizon). The
        result is statistically clean and economically meaningful
        per unit, but it does not explain large fractions of
        return variance.
  \item I do not have a structural model. The narrative
        explanation of the U-shape is a story, not a derivation.
\end{enumerate}

The questions I'd most value your read on:

\begin{itemize}
  \item Does the U-shape result, by itself, count as a
        contribution? It overturns a piece of folk wisdom in the
        innovation literature (the ``pioneers pay; fast followers
        win'' story) by saying that both pioneers and templated
        followers do better than the middle. If you have seen
        this in print before, please point me; I have not.
  \item Is the right venue management/strategy
        (\emph{Strategic Management Journal},
        \emph{Organization Science}), economics
        (\emph{NBER Working Paper}, \emph{AER P\&P}), or
        information science (\emph{Management Science},
        \emph{ICWSM} for the methodological piece)?
  \item Should I split the methodological contribution
        (KL-surprise on incentivised text) from the empirical
        finding (U-shape, additivity-with-patents, multi-year
        return cycle), or are they better together?
  \item What's the right second study? My current short list:
        a 10-K replication, a Kickstarter-vs-equity-outcome
        comparison, or a more careful entity-resolution pass
        with Compustat to enlarge the matched panel and add risk-
        adjusted return measures (Fama--French factors) to the
        stock-return regression.
\end{itemize}

The whole thing is in
\href{https://github.com/ericsilver/novelty}{\texttt{github.com/ericsilver/novelty}};
the longer write-up is in \texttt{paper/main.pdf}. About thirty
pages. I'd rather not lose another decade between thinking about
this and finishing it.

\medskip
With thanks,

\medskip
\noindent Eric

\end{document}
